\chapter{The Zeros of \texorpdfstring{$\zeta$}{zeta}}

\begin{center}\itshape
One now finds indeed approximately this number of real roots within these limits, and it is very probable that all roots are real.\\
\upshape --- Riemann, \emph{On the number of primes less than a given magnitude} (1859)
\end{center}

Riemann's memoir is eight pages long, and it is one of the strangest documents in mathematics. In those pages he extends $\zeta$ to the whole plane, proves the
functional equation, and then, almost in passing, makes three claims about the zeros in the strip $0<\operatorname{Re}s<1$. There are infinitely many of them. The number
with ordinate between $0$ and $T$ is about $\frac T{2\pi}\log\frac T{2\pi}-\frac T{2\pi}$. And the function $\xi(s)$, which vanishes exactly at those zeros, is a product over them, like a
polynomial over its roots. Riemann gave no proofs; he wrote that he had put the question of the reality of the roots aside ``after some fleeting vain attempts''.

The second and third claims took thirty-four and forty-six years to prove. Hadamard proved the product in 1893, inventing for the purpose the theory of entire functions
of finite order; three years later he used it, as de la Vallée Poussin did, to prove the prime number theorem. Von Mangoldt proved the counting formula in 1905. The
first claim, the one about the real roots, is the Riemann hypothesis, and it is still open.

This chapter proves the two theorems that are proved, and one more: the zeros stay away from the line $\operatorname{Re}s=1$ by a margin, which Chapter~7 turns into
an error term for the prime number theorem. The tools are four classical pieces of complex analysis that Volume~4 did not need: Stirling's formula off the real line,
Jensen's formula, which counts zeros by averaging $\log|f|$ on a circle, the Borel--Carathéodory lemma, which bounds a function by its real part, and the argument principle
of Volume~4, Theorem~9.37.

\begin{think}
Figure~\ref{fig:nt} shows that $N(T)$, the number of zeros with ordinate in $]0,T]$, follows Riemann's curve closely. Before reading on: why would the number of zeros
grow like $T\log T$, faster than linearly? Where in the functional equation could a factor $\log T$ come from?
\end{think}

\begin{figure}[ht]
\centering
\fig{ch06-nt}
\caption{The staircase $N(T)$ (blue) for $0\le T\le100$, with $29$ steps, one at each zero $\frac12+i\gamma$: the first at $\gamma=14.13$, the last before $100$ at $\gamma=98.83$. The smooth
curve (orange) is Riemann's main term $\frac T{2\pi}\log\frac T{2\pi e}+\frac78$. The difference $S(T)$ stays small: $|S(T)|<0.83$ for $1\le T\le100$.}
\label{fig:nt}
\end{figure}

\section{Stirling's formula in the complex plane}

The functional equation multiplies $\zeta$ by $\pi^{-s/2}\Gamma(\frac s2)$, so every question about $\zeta$ far up the critical strip becomes, in part, a question about
$\Gamma(\sigma+it)$ for large $t$. Volume~4 proved Stirling's formula for real arguments only. Here is the complex version; the trick is to write down the right
holomorphic function and let the identity theorem do the rest.

For $s$ in the slit plane $\C^-=\C\setminus\left]-\infty,0\right]$ let $\log s$ be the principal logarithm. Let $P(u)=\frac12(\{u\}^2-\{u\})$: it is continuous, $1$-periodic, with
$-\frac18\le P\le0$ and $P'(u)=\{u\}-\frac12$ away from the integers. Define, for $s\in\C^-$,
\[
\begin{gathered}J(s)=-\int_0^\infty\frac{P(u)}{(u+s)^2}\dd u,\qq F(s)=\Bigl(s-\frac12\Bigr)\log s-s+\frac12\log2\pi+J(s).\end{gathered}
\]

\begin{theorem}[Stirling's formula]\label{thm:stirlingc}
\begin{enumerate}[label=(\alph*)]
\item $F$ is holomorphic on $\C^-$ and $e^{F}=\Gamma$ there. So $F$ is a holomorphic logarithm of $\Gamma$ on $\C^-$, real on $\left]0,\infty\right[$, which we write $\log\Gamma$.
\item Let $0<\delta\le\frac\pi2$. For $|\arg s|\le\pi-\delta$, $|J(s)|\le\frac{C_\delta}{|s|}$ and $|J'(s)|\le\frac{C_\delta}{|s|^2}$, with $C_\delta=(1-\cos\delta)^{-3/2}$.
\item Consequently, for $|\arg s|\le\pi-\delta$ and $|s|\ge1$,
\[
\begin{aligned}\log\Gamma(s)&=\Bigl(s-\frac12\Bigr)\log s-s+\frac12\log2\pi\eqbr+O_\delta\Bigl(\frac1{|s|}\Bigr),\\\frac{\Gamma'}\Gamma(s)&=\log s+O_\delta\Bigl(\frac1{|s|}\Bigr).\end{aligned}
\]
\end{enumerate}
\end{theorem}
\begin{proof}
\emph{Holomorphy.} Let $K\subset\C^-$ be compact, at distance $d>0$ from $\left]-\infty,0\right]$, with $|s|\le S$ on $K$. For $u\ge0$ and $s\in K$, $|u+s|\ge d$, and $|u+s|\ge u-S\ge\frac u2$ when
$u\ge2S$; so $|u+s|\ge c(1+u)$ for some $c>0$. The integrand of $J$ is continuous in $(u,s)$, holomorphic in $s$, and bounded by $\frac1{8c^2(1+u)^2}$, which is integrable:
Volume~4, Theorem~9.15(b) shows that $J$ is holomorphic on $\C^-$, and so is $F$.

\emph{The functional equation.} We show $F(s+1)=F(s)+\log s$. Substituting $v=u+1$, $J(s+1)=-\int_1^\infty\frac{P(v)}{(v+s)^2}\dd v$, so, integrating by parts with $P(0)=P(1)=0$,
\[
\begin{aligned}J(s+1)-J(s)&=\int_0^1\frac{P(u)}{(u+s)^2}\dd u\eqbr=\int_0^1\frac{u-\frac12}{u+s}\dd u\\&=1-\Bigl(s+\frac12\Bigr)\eqbr\quad\times\bigl(\log(s+1)-\log s\bigr),\end{aligned}
\]
since $\frac{u-1/2}{u+s}=1-\frac{s+1/2}{u+s}$ and $\log(u+s)$ is a primitive of $\frac1{u+s}$ along $[s,s+1]\subset\C^-$. Adding $(s+\frac12)\log(s+1)-(s-\frac12)\log s-1$ gives
$F(s+1)-F(s)=\log s$.

\emph{Identification.} For real $x>0$, $H(x)=F(x)-\log\Gamma(x)$ is real, and $H(x+1)=H(x)$ by the functional equations of $F$ and $\Gamma$. As $x\to\infty$, $J(x)\to0$
(by (b) below, or directly: $|J(x)|\le\frac18\int_0^\infty\frac{\mathrm du}{(u+x)^2}=\frac1{8x}$) and $\log\Gamma(x)-(x-\frac12)\log x+x-\frac12\log2\pi\to0$ by Volume~4, Theorem~5.8. So the
periodic function $H$ tends to $0$, hence vanishes: $e^F=\Gamma$ on $\left]0,\infty\right[$. Both sides are holomorphic on the connected set $\C^-$ (Volume~4, Theorem~9.36(b)), so they
agree there (Volume~4, Theorem~9.18).

(b) Write $s=re^{i\phi}$ with $|\phi|\le\pi-\delta$, and $c=\cos\phi\ge-\cos\delta$. For $u\ge0$, $|u+s|^2=u^2+r^2+2urc\ge(1-\cos\delta)(u^2+r^2)$: this is clear if $c\ge0$ (as $\cos\delta\ge0$), and if
$c<0$ it follows from $2ur\le u^2+r^2$. Hence, with $\kappa=1-\cos\delta$,
\[
\begin{aligned}|J(s)|&\le\frac18\int_0^\infty\frac{\mathrm du}{\kappa(u^2+r^2)}=\frac\pi{16\kappa r},\\|J'(s)|&=\Bigl|2\int_0^\infty\frac{P(u)}{(u+s)^3}\dd u\Bigr|\eqbr\le\frac14\int_0^\infty\frac{\mathrm du}{\kappa^{3/2}(u^2+r^2)^{3/2}}=\frac1{4\kappa^{3/2}r^2},\end{aligned}
\]
using $\int_0^\infty\frac{\mathrm du}{(u^2+r^2)^{3/2}}=\frac1{r^2}$ (substitute $u=r\tan\theta$). Both are at most $\frac{C_\delta}{|s|}$ and $\frac{C_\delta}{|s|^2}$.

(c) The first formula is (a) with (b). Differentiating $F$: $\frac{\Gamma'}\Gamma(s)=F'(s)=\log s+\frac{s-1/2}s-1+J'(s)=\log s-\frac1{2s}+J'(s)$.
\end{proof}

On a vertical line, Stirling's formula becomes very concrete. Let $\sigma$ be fixed and $t\to+\infty$. Then $\log(\sigma+it)=\log t+i\frac\pi2-\frac{i\sigma}t+O(\frac1{t^2})$, and
multiplying out,
\begin{equation}\label{eq:stirlingvert}
\begin{aligned}\log\Gamma(\sigma+it)&=\Bigl(\sigma-\frac12\Bigr)\log t-\frac{\pi t}2+\frac12\log2\pi\\&\quad+i\Bigl(t\log t-t+\frac\pi2\Bigl(\sigma-\frac12\Bigr)\Bigr)\eqbr+O\Bigl(\frac1t\Bigr).\end{aligned}
\end{equation}
So $|\Gamma(\sigma+it)|\sim\sqrt{2\pi}\,t^{\sigma-1/2}e^{-\pi t/2}$: the exponential decay of Volume~4, Theorem~9.36(d), now on every vertical line, and with the polynomial factor that
decides everything in this chapter.

\section{The function \texorpdfstring{$\xi$}{xi}}

Riemann's $\xi$ packs $\zeta$, the Gamma factor and the pole into one entire function. Volume~4, Theorem~10.8 shows that $\Lambda(s)=\pi^{-s/2}\Gamma(\frac s2)\zeta(s)$ satisfies
$\Lambda(s)=-\frac1s-\frac1{1-s}+I(s)$ with $I(s)=\int_1^\infty\omega(t)(t^{s/2}+t^{(1-s)/2})\frac{\mathrm dt}t$ entire and $I(1-s)=I(s)$. (In this volume $\Lambda$ is von Mangoldt's function, so we
use this notation only here.) Let
\[
\xi(s)=\frac12s(s-1)\pi^{-s/2}\Gamma\Bigl(\frac s2\Bigr)\zeta(s).
\]

\begin{proposition}\label{prop:xi}
\begin{enumerate}[label=(\alph*)]
\item $\xi$ is entire, $\xi(s)=\xi(1-s)$, $\xi(\overline s)=\overline{\xi(s)}$, and $\xi(0)=\xi(1)=\frac12$.
\item The zeros of $\xi$ are exactly the zeros $\rho=\beta+i\gamma$ of $\zeta$ with $0<\beta<1$, with the same multiplicities. They are symmetric with respect to the real axis and to
the line $\operatorname{Re}s=\frac12$, and none is real.
\item There is a constant $C$ such that $|\xi(s)|\le e^{C|s|\log|s|}$ for $|s|\ge2$.
\end{enumerate}
\end{proposition}
\begin{proof}
(a) $\xi(s)=\frac12s(s-1)\Lambda(s)=\frac12s(s-1)I(s)+\frac12\bigl(-(s-1)+s\bigr)=\frac12+\frac12s(s-1)I(s)$, which is entire and symmetric under $s\mapsto1-s$, and equals $\frac12$ at $0$ and $1$.
Conjugation symmetry holds for $\operatorname{Re}s>1$, where every factor is real on the real axis, hence everywhere by the identity theorem.

(b) For $\operatorname{Re}s>1$: $\zeta(s)\ne0$ (Euler product), $\Gamma(\frac s2)\ne0$ (Volume~4, Theorem~9.36(c)), so $\xi(s)\ne0$. By symmetry $\xi\ne0$ on $\operatorname{Re}s<0$. On $\operatorname{Re}s=1$:
$\zeta(1+it)\ne0$ for $t\ne0$ (Theorem~\ref{thm:line}) and $\xi(1)=\frac12$; by symmetry nothing on $\operatorname{Re}s=0$. In the open strip $s(s-1)\pi^{-s/2}\Gamma(\frac s2)$ is holomorphic
and non-zero, so $\xi$ and $\zeta$ have the same zeros there, with the same orders. Finally $\zeta(\sigma)<0$ for $0<\sigma<1$ (Volume~4, Exercise~10.3), so no zero is real.

(c) If $\sigma<\frac12$ then $|s|\le|1-s|\le|s|+1$, so by symmetry it suffices to treat $\sigma\ge\frac12$, $|s|\ge2$, where $|s-1|\ge1$. The proof of Volume~4, Proposition~10.6 gives, for $\sigma>0$,
$|\zeta(s)-\frac1{s-1}|\le|s|\sum_nn^{-\sigma-1}\le|s|(1+\frac1\sigma)$, so $|\zeta(s)|\le1+3|s|$ when $\sigma\ge\frac12$ and $|s-1|\ge1$. Next $|\pi^{-s/2}|\le1$, and
$|\Gamma(\frac s2)|\le\Gamma(\frac\sigma2)$ (Volume~4, Theorem~9.36(a)). For $x\ge1$, $t^{x-1}e^{-t/2}\le(2x)^{x-1}$ for all $t>0$ (the maximum is at $t=2(x-1)$), so $\Gamma(x)\le(2x)^{x-1}\int_0^\infty e^{-t/2}\dd t\le(2x)^x$; for
$\frac14\le x\le1$, $\Gamma(x)\le\Gamma(\frac14)<4$ because $\Gamma$ is convex and $\Gamma(1)=1$. Altogether $|\xi(s)|\le|s|^2(1+3|s|)\bigl(4+(2|s|)^{|s|}\bigr)$ there, which is
at most $e^{C|s|\log|s|}$ for $|s|\ge2$.
\end{proof}

The bound (c) is sharp in shape: for real $\sigma\to\infty$, $\xi(\sigma)\ge\frac12\sigma(\sigma-1)\pi^{-\sigma/2}\Gamma(\frac\sigma2)$, which grows like $e^{\frac12\sigma\log\sigma(1+o(1))}$ by Stirling. We say that
$\xi$ has \emph{order $1$}. That single number is what makes the next three sections work.

\section{Counting zeros with a circle}

How can one count the zeros of a function without finding them? In 1899 the Danish engineer Jensen, who worked for the Copenhagen telephone company and did
mathematics in his spare time, gave the answer: average $\log|f|$ over a circle.

\begin{theorem}[Jensen's formula]\label{thm:jensen}
Let $f$ be holomorphic on an open set containing the closed disc $|z|\le R$, with $f(0)\ne0$, and let $z_1,\dots,z_n$ be its zeros in $|z|<R$, repeated according to multiplicity.
If $f$ has no zero on $|z|=R$, then
\[
\begin{aligned}\log|f(0)|&+\sum_{k=1}^n\log\frac R{|z_k|}\eqbr=\frac1{2\pi}\int_0^{2\pi}\log|f(Re^{i\theta})|\dd\theta.\end{aligned}
\]
\end{theorem}
\begin{proof}
For $|a|<R$, the \emph{Blaschke factor} $B_a(z)=\frac{R(z-a)}{R^2-\overline az}$ is holomorphic on $|z|<\frac{R^2}{|a|}$, a disc larger than $|z|\le R$, vanishes only at $a$, and has modulus $1$ on
$|z|=R$: there $R^2=z\overline z$, so $|R^2-\overline az|=|z||\overline z-\overline a|=R|z-a|$. Let $g=f/\prod_kB_{z_k}$. It is holomorphic on an open disc $D$ slightly larger than $|z|\le R$
(the zeros cancel, with their multiplicities), has no zeros on the closed disc, hence none on a slightly smaller open disc $D'\supset\{|z|\le R\}$. On the convex set $D'$, $\frac{g'}g$ has a
primitive $\ell$ (Volume~4, Theorem~9.9), and $(ge^{-\ell})'=0$, so after adding a constant $e^\ell=g$: $\ell$ is a holomorphic logarithm of $g$. The mean value property (Cauchy's formula at
the centre) gives $\ell(0)=\frac1{2\pi}\int_0^{2\pi}\ell(Re^{i\theta})\dd\theta$; take real parts: $\log|g(0)|=\frac1{2\pi}\int\log|g(Re^{i\theta})|\dd\theta$. Finally $|g|=|f|$ on the circle, and
$|g(0)|=|f(0)|\prod_k\frac R{|z_k|}$ since $|B_a(0)|=\frac{|a|}R$.
\end{proof}

We use it as an inequality, which needs no condition on the circle.

\begin{corollary}\label{cor:jensen}
Let $f$ be holomorphic on an open set containing $|z|\le R$, with $f(0)\ne0$ and $|f|\le M$ on $|z|=R$. For $0<r<R$, the number $n(r)$ of zeros of $f$ in $|z|\le r$ satisfies
\[
n(r)\log\frac Rr\le\log\frac M{|f(0)|}.
\]
\end{corollary}
\begin{proof}
The zeros in the compact disc are finitely many, so there are radii $R'\in\left]r,R\right[$ arbitrarily close to $R$ with no zero on $|z|=R'$. In Jensen's formula for $R'$, each of the $n(r)$
zeros in $|z|\le r$ contributes at least $\log\frac{R'}r$, the others contribute non-negative terms, and the right side is at most $\log\max_{|z|=R'}|f|\le\log M$ by the maximum
principle (Volume~4, Theorem~9.19). Let $R'\to R$.
\end{proof}

\begin{beauty}
Jensen's formula is a conservation law. The average of $\log|f|$ on a circle exceeds its value at the centre by exactly the amount that the zeros inside ``cost'', each zero
at distance $|z_k|$ costing $\log\frac R{|z_k|}$. A function that is not allowed to grow fast is not allowed many zeros.
\end{beauty}

Applied to $\xi$ with $|\xi(0)|=\frac12$ and Proposition~\ref{prop:xi}(c) on the circle of radius $2r$, Corollary~\ref{cor:jensen} gives
\begin{equation}\label{eq:nr}
\#\{\rho:|\rho|\le r\}\le\frac{\log(2e^{2Cr\log2r})}{\log2}=O(r\log r).
\end{equation}
Summing by parts, $\sum_\rho|\rho|^{-1-\varepsilon}<\infty$ for every $\varepsilon>0$; in particular $\sum_\rho|\rho|^{-2}<\infty$. (Let $r_0>0$ be the smallest modulus of a zero, which exists since $\xi(0)\ne0$ and the zeros are isolated; then
$\sum_{|\rho|\le R}|\rho|^{-2}=\int_{r_0^-}^R\frac{\mathrm dn(r)}{r^2}=\frac{n(R)}{R^2}+2\int_{r_0}^R\frac{n(r)}{r^3}\dd r$, which stays bounded.) Exercise~\ref{ex:zerosums} shows that
$\sum|\rho|^{-1}$ diverges: there are many zeros, but not too many.

\section{Hadamard's product}

A polynomial is determined, up to a constant, by its roots. An entire function is not: $e^z$ has no zeros at all. Hadamard's discovery is that for functions that do not
grow too fast, the ambiguity is only an exponential factor of controlled degree. We prove it for $\xi$, in the form of a partial fraction expansion of $\frac{\xi'}\xi$, from
which the product follows. The tool is a lemma that bounds a holomorphic function by its real part alone.

\begin{lemma}[Borel--Carathéodory]\label{lem:borel}
Let $h$ be holomorphic on an open set containing $|z|\le R$, with $h(0)=0$ and $\operatorname{Re}h\le A$ on $|z|=R$. Then for $|z|\le r<R$ and $k\ge1$,
\[
|h^{(k)}(z)|\le\frac{2A\,k!\,R}{(R-r)^{k+1}}.
\]
\end{lemma}
\begin{proof}
$h$ is holomorphic on a disc $|z|<R_1$ with $R_1>R$, so $h(z)=\sum_{n\ge1}a_nz^n$ there (Volume~4, Theorem~9.11), uniformly on $|z|=R$. Integrating term by term against $e^{-in\theta}$ on
the circle,
\[
\begin{gathered}a_nR^n=\frac1{2\pi}\int_0^{2\pi}h(Re^{i\theta})e^{-in\theta}\dd\theta,\qq0=\frac1{2\pi}\int_0^{2\pi}\overline{h(Re^{i\theta})}e^{-in\theta}\dd\theta\end{gathered}
\]
for $n\ge1$, the second because its conjugate is the coefficient of $z^{-n}$, which is $0$. Adding, $a_nR^n=\frac1\pi\int_0^{2\pi}\operatorname{Re}h(Re^{i\theta})e^{-in\theta}\dd\theta$. Since
$\int_0^{2\pi}e^{-in\theta}\dd\theta=0$, also $a_nR^n=-\frac1\pi\int_0^{2\pi}\bigl(A-\operatorname{Re}h(Re^{i\theta})\bigr)e^{-in\theta}\dd\theta$. Now $A-\operatorname{Re}h\ge0$ on the circle, and its integral is
$2\pi A$, since $\int\operatorname{Re}h=2\pi\operatorname{Re}h(0)=0$. So $|a_n|R^n\le2A$, and
\[
\begin{aligned}|h^{(k)}(z)|&\le\sum_{n\ge k}n(n-1)\cdots(n-k+1)\frac{2A}{R^n}r^{n-k}\eqbr=\frac{2A}{R^k}\cdot\frac{k!}{(1-r/R)^{k+1}},\end{aligned}
\]
using $\sum_{n\ge k}\binom nkx^{n-k}=(1-x)^{-k-1}$. This is the claimed bound.
\end{proof}

The point is that only an \emph{upper} bound on $\operatorname{Re}h$ is needed. For $h=\log\frac{f}{f(0)}$ that is an upper bound on $|f|$, which is what growth estimates give.

\begin{lemma}\label{lem:localzeros}
Let $f$ be holomorphic on an open set containing $|z|\le R$, with $f(0)\ne0$ and $|f(z)|\le M|f(0)|$ for $|z|\le R$, where $M\ge2$. Let $\mathcal Z$ be the set of zeros of $f$ in $|z|\le\frac R2$,
with multiplicity. Then $\#\mathcal Z\le\frac{\log M}{\log2}$, and for $|z|\le\frac R4$, $z\notin\mathcal Z$,
\[
\begin{gathered}\Bigl|\frac{f'}f(z)-\sum_{\rho\in\mathcal Z}\frac1{z-\rho}\Bigr|\le\frac{17\log M}R,\\\Bigl|\Bigl(\frac{f'}f\Bigr)'(z)+\sum_{\rho\in\mathcal Z}\frac1{(z-\rho)^2}\Bigr|\le\frac{129\log M}{R^2}.\end{gathered}
\]
\end{lemma}
\begin{proof}
The count is Corollary~\ref{cor:jensen} with $r=\frac R2$. Let $g(z)=f(z)\prod_{\rho\in\mathcal Z}\frac{R^2-\overline\rho z}{R(z-\rho)}$. Each factor has modulus $1$ on $|z|=R$ (as for the Blaschke factors), and its
pole at $\rho$ is cancelled by a zero of $f$; its zero $\frac{R^2}{\overline\rho}$ has modulus at least $2R$. So $g$ is holomorphic near $|z|\le R$, $|g|=|f|\le M|f(0)|$ on $|z|=R$, hence on the
disc by the maximum principle, and $|g(0)|=|f(0)|\prod\frac R{|\rho|}\ge|f(0)|$. Also $g$ has no zero in $|z|\le\frac R2$, so, as in the proof of Jensen's formula, $h=\log\frac g{g(0)}$ is
holomorphic near $|z|\le\frac R2$, with $h(0)=0$ and $\operatorname{Re}h=\log\frac{|g|}{|g(0)|}\le\log M$. Lemma~\ref{lem:borel} with radius $\frac R2$ and $r=\frac R4$ gives $|h'|\le\frac{16\log M}R$ and
$|h''|\le\frac{128\log M}{R^2}$ on $|z|\le\frac R4$. Now
\[
h'=\frac{f'}f-\sum_{\rho\in\mathcal Z}\frac1{z-\rho}-\sum_{\rho\in\mathcal Z}\frac{\overline\rho}{R^2-\overline\rho z},
\]
and for $|z|\le\frac R4$, $|\rho|\le\frac R2$: $\bigl|\frac{\overline\rho}{R^2-\overline\rho z}\bigr|\le\frac{R/2}{R^2-R^2/8}=\frac4{7R}$. With at most $\frac{\log M}{\log2}$ such terms, the last sum is at most
$\frac{4}{7\log2}\frac{\log M}R<\frac{\log M}R$. Differentiating once more, the corresponding terms are at most $\frac{16}{49R^2}$ each, and $\frac{16}{49\log2}<1$.
\end{proof}

\begin{theorem}[Hadamard, 1893]\label{thm:hadamard}
Let $\rho$ run through the zeros of $\xi$, with multiplicity, and let $B=\frac{\xi'(0)}{\xi(0)}$. Then, for every $s$ that is not a zero,
\[
\frac{\xi'(s)}{\xi(s)}=B+\sum_\rho\Bigl(\frac1{s-\rho}+\frac1\rho\Bigr),
\]
the series converging absolutely and uniformly on compact sets avoiding the zeros. Equivalently,
\[
\xi(s)=\frac12e^{Bs}\prod_\rho\Bigl(1-\frac s\rho\Bigr)e^{s/\rho}.
\]
\end{theorem}
\begin{proof}
\emph{The second derivative.} Fix $s$, not a zero, and let $R\ge4|s|+8$. On $|z|\le R$, $|\xi(z)|\le e^{CR\log R}=2e^{CR\log R}|\xi(0)|$, so Lemma~\ref{lem:localzeros} applies with $\log M=O(R\log R)$:
\[
\Bigl(\frac{\xi'}\xi\Bigr)'(s)=-\sum_{|\rho|\le R/2}\frac1{(s-\rho)^2}+O\Bigl(\frac{\log R}R\Bigr).
\]
As $R\to\infty$ the error tends to $0$, and the sum tends to the full series, which converges absolutely: for $|\rho|\ge2|s|$, $|s-\rho|\ge\frac{|\rho|}2$, and $\sum|\rho|^{-2}<\infty$. So
$\bigl(\frac{\xi'}\xi\bigr)'=-\sum_\rho(s-\rho)^{-2}$.

\emph{Integrating once.} Let $\Phi(s)=\sum_\rho\bigl(\frac1{s-\rho}+\frac1\rho\bigr)=\sum_\rho\frac s{\rho(s-\rho)}$. On a compact set $K$ avoiding the zeros, with $|s|\le S$, the terms with $|\rho|\ge2S$ are at
most $\frac{2S}{|\rho|^2}$: the series converges uniformly on $K$, so $\Phi$ is holomorphic off the zeros, with $\Phi'=-\sum(s-\rho)^{-2}$ (Volume~4, Theorem~9.15(a)). Hence
$\frac{\xi'}\xi-\Phi$ has derivative $0$ on the complement of the zeros, which is connected (the zeros are isolated), so it is a constant, its value $B$ at $s=0$.

\emph{The product.} For $|w|\le\frac12$, the principal value $\ell(w)=\log(1-w)+w=-\sum_{k\ge2}\frac{w^k}k$ satisfies $|\ell(w)|\le|w|^2$. For $|s|\le S$ and $|\rho|\ge2S$, the factor
$(1-\frac s\rho)e^{s/\rho}$ is therefore $e^{\ell(s/\rho)}$ with $\sum|\ell(s/\rho)|\le S^2\sum|\rho|^{-2}$: the product converges uniformly on $|s|\le S$ to a holomorphic function $E(s)$ whose zeros are
the $\rho$, with $\frac{E'}E=\Phi$ (differentiate the uniformly convergent sum of logarithms, and the finitely many other factors directly). Then $\frac\xi E$ is entire without zeros and its
logarithmic derivative is $B$, so $\xi=\xi(0)e^{Bs}E$.
\end{proof}

\begin{corollary}\label{cor:infinitezeros}
$\zeta$ has infinitely many zeros in the strip $0<\operatorname{Re}s<1$.
\end{corollary}
\begin{proof}
If there were finitely many, $E$ would be a polynomial times $e^{cs}$, and $|\xi(\sigma)|\le e^{C'\sigma}$ for real $\sigma\ge1$. But $\xi(\sigma)\ge\frac12\sigma(\sigma-1)\pi^{-\sigma/2}\Gamma(\frac\sigma2)$ grows faster
than any exponential, by Stirling's formula.
\end{proof}

Taking the logarithmic derivative of the definition of $\xi$ gives the form in which the theorem is used. With $\psi_\Gamma=\frac{\Gamma'}\Gamma$ (the letter $\psi$ is taken),
\begin{equation}\label{eq:zetapf}
\begin{aligned}\frac{\zeta'}\zeta(s)&=B+\sum_\rho\Bigl(\frac1{s-\rho}+\frac1\rho\Bigr)-\frac1s-\frac1{s-1}\eqbr+\frac12\log\pi-\frac12\psi_\Gamma\Bigl(\frac s2\Bigr).\end{aligned}
\end{equation}
The pole of $\zeta$ appears as $-\frac1{s-1}$, the trivial zeros come from the poles of $\psi_\Gamma(\frac s2)$ at $0,-2,-4,\dots$ (the one at $0$ cancels $-\frac1s$), and every non-trivial zero contributes
a simple pole. Exercises~\ref{ex:digammahalf} and~\ref{ex:bsum} show that $B=\frac12\log4\pi-1-\frac\gamma2=-0.0230957\ldots=-\sum_\rho\operatorname{Re}\frac1\rho$.

\section{The Riemann--von Mangoldt formula}

Order the ordinates $\gamma>0$ of the zeros, and let $N(T)$ be the number of zeros $\rho=\beta+i\gamma$ with $0<\gamma\le T$, counted with multiplicity. Following Volume~4, Chapter~10, let
\[
\vartheta(T)=\operatorname{Im}\log\Gamma\Bigl(\frac14+\frac{iT}2\Bigr)-\frac T2\log\pi,
\]
with the logarithm of Theorem~\ref{thm:stirlingc}. By \eqref{eq:stirlingvert} with $\sigma=\frac14$ and $t=\frac T2$,
\begin{equation}\label{eq:theta}
\vartheta(T)=\frac T2\log\frac T{2\pi}-\frac T2-\frac\pi8+O\Bigl(\frac1T\Bigr).
\end{equation}

\begin{theorem}[Riemann 1859, von Mangoldt 1905]\label{thm:rvm}
For $T\ge3$ not the ordinate of a zero,
\[
\begin{gathered}N(T)=\frac{\vartheta(T)}\pi+1+S(T),\qq S(T)=\frac1\pi\operatorname{Im}\int_L\frac{\zeta'(s)}{\zeta(s)}\dd s,\end{gathered}
\]
where $L$ is the path from $2$ to $2+iT$ to $\frac12+iT$, and $S(T)=O(\log T)$. Hence
\[
N(T)=\frac T{2\pi}\log\frac T{2\pi e}+\frac78+O(\log T).
\]
\end{theorem}

The constant $\frac78$ is no typo: Riemann's main term gets an $\frac78$ from the Gamma factor and the pole. The real content is that the error is only $O(\log T)$, while the main term
grows like $T\log T$.

\begin{proof}
\emph{Counting with the argument principle.} Let $D$ be the open rectangle $-1<\sigma<2$, $0<t<T$. The zeros of $\xi$ in $D$ are those of $\zeta$ with $0<\gamma<T$ (none is real), and none
lies on $\partial D$. By Volume~4, Theorem~9.37, $N(T)=\frac1{2\pi i}\int_{\partial D}\frac{\xi'}\xi$, which is real, so $2\pi N(T)=\operatorname{Im}\int_{\partial D}\frac{\xi'}\xi$. On the bottom side $\xi$ is real and
positive ($\xi(0)=\frac12$, and no real zeros), so $\operatorname{Im}\frac{\xi'}\xi=0$ there. The map $s\mapsto1-\overline s$ exchanges the right half of the rest of $\partial D$, which is $L$, with the left
half, reversing the direction; since $\xi(1-\overline s)=\overline{\xi(s)}$, $\operatorname{Im}\frac{\xi'}\xi$ integrates to the same value on both. So $\pi N(T)=\operatorname{Im}\int_L\frac{\xi'}\xi$.

\emph{The four factors.} $\frac{\xi'}\xi=\frac1s+\frac1{s-1}-\frac12\log\pi+\frac12\psi_\Gamma(\frac s2)+\frac{\zeta'}\zeta$. Along $L$:
\begin{itemize}
\item $\operatorname{Im}\int_L\bigl(\frac1s+\frac1{s-1}\bigr)$ is the change of a continuous argument of $s(s-1)$. It starts at $2>0$, ends at $(\frac12+iT)(-\frac12+iT)=-T^2-\frac14<0$, and in between
$\operatorname{Im}s(s-1)=t(2\sigma-1)>0$: the change is exactly $\pi$.
\item $\operatorname{Im}\int_L(-\frac12\log\pi)\dd s=-\frac12\log\pi\cdot\operatorname{Im}(\frac12+iT-2)=-\frac T2\log\pi$.
\item $\int_L\frac12\psi_\Gamma(\frac s2)\dd s=\log\Gamma(\frac14+\frac{iT}2)-\log\Gamma(1)$, whose imaginary part is $\operatorname{Im}\log\Gamma(\frac14+\frac{iT}2)$, since $L$ stays in $\C^-$.
\end{itemize}
Adding, $\pi N(T)=\pi+\vartheta(T)+\pi S(T)$.

\emph{The bound on $S(T)$.} On the vertical side, $\operatorname{Re}\zeta(2+it)\ge1-\sum_{n\ge2}n^{-2}=2-\zeta(2)>0.35$, so a continuous argument of $\zeta$ stays in $\left]-\frac\pi2,\frac\pi2\right[$ and changes by
less than $\pi$. On the horizontal side, suppose $\operatorname{Re}\zeta(\sigma+iT)$ vanishes at $q$ points of $[\frac12,2]$. They cut the segment into at most $q+1$ pieces, on each of which
$\operatorname{Re}\zeta$ has constant sign, so the argument changes by less than $\pi$ on each: $\pi|S(T)|<\pi+(q+1)\pi$.

To bound $q$, let $k(z)=\frac12\bigl(\zeta(z+iT)+\zeta(z-iT)\bigr)$, holomorphic on $|z-2|\le\frac74$ (the poles of the two terms are at $1\mp iT$, far away). For real $z=\sigma$,
$\zeta(\sigma-iT)=\overline{\zeta(\sigma+iT)}$, so $k(\sigma)=\operatorname{Re}\zeta(\sigma+iT)$, and $q$ is at most the number of zeros of $k$ in $|z-2|\le\frac32$. On $|z-2|\le\frac74$, $\operatorname{Re}(z\pm iT)\ge\frac14$ and
$|z\pm iT|\le T+4$, so $|k(z)|\le1+5(T+4)$ by the bound of Proposition~\ref{prop:xi}(c) (whose proof gives $|\zeta(s)|\le\frac1{|s-1|}+5|s|$ for $\sigma\ge\frac14$). And $k(2)>0.35$. Corollary~\ref{cor:jensen}
gives $q\le\frac{\log(3(5T+21))}{\log(7/6)}=O(\log T)$.

Finally, \eqref{eq:theta} turns $\frac{\vartheta(T)}\pi+1$ into $\frac T{2\pi}\log\frac T{2\pi}-\frac T{2\pi}-\frac18+1+O(\frac1T)$.
\end{proof}

\begin{corollary}\label{cor:nunit}
For $T\ge3$, the number of zeros with $T<\gamma\le T+1$ is $O(\log T)$. For every $s=\sigma+it$ with $|t|\ge2$, $\sum_\rho\frac1{1+(t-\gamma)^2}=O(\log|t|)$.
\end{corollary}
\begin{proof}
$N(T+1)-N(T)=\frac{\vartheta(T+1)-\vartheta(T)}\pi+O(\log T)$, and $\vartheta'(u)=\frac12\log\frac u{2\pi}+O(\frac1u)$ by Theorem~\ref{thm:stirlingc}(c) (differentiate under the imaginary part). For the
sum, the zeros are symmetric about the real axis, so it is at most twice the sum over $\gamma>0$; group the $\gamma$ with $k\le|\gamma-t|<k+1$, $k\ge0$: there are $O(\log(|t|+k+2))$ of them, contributing
$O\bigl(\frac{\log(|t|+k+2)}{1+k^2}\bigr)$, and $\sum_k\frac{\log(|t|+k+2)}{1+k^2}=O(\log|t|)$.
\end{proof}

\begin{remark}
On the critical line $\xi(\frac12+it)$ is real, and its sign changes are zeros on the line. Volume~4 found them through Hardy's function $Z(t)=e^{i\vartheta(t)}\zeta(\frac12+it)$. Theorem~\ref{thm:rvm}
gives the total count in the strip, so if a computer finds $N(T)$ sign changes of $Z$ on $]0,T]$, then \emph{all} zeros up to height $T$ are on the line and simple. Making this rigorous
needs care with $S(T)$ at the endpoint; Turing found in 1953 how, using that $S(t)$ has mean $0$ (Littlewood's $\int_0^TS=O(\log T)$). By this method the first $10^{13}$ zeros have been
checked. A numerical calculation does it for the first $79$ zeros, up to $T=200$.
\end{remark}

\section{Zeros near a point, and away from the line}

The partial fraction \eqref{eq:zetapf} is an identity between two divergent-looking objects. Taking differences makes it local: near a point $s$, $\frac{\zeta'}\zeta$ is the
sum of $\frac1{s-\rho}$ over the zeros close to $s$, up to $O(\log|t|)$.

\begin{proposition}\label{prop:nearzeros}
Uniformly for $-1\le\sigma\le2$ and $|t|\ge2$,
\[
\frac{\zeta'}\zeta(\sigma+it)=\sum_{|\gamma-t|\le1}\frac1{s-\rho}+O(\log|t|).
\]
\end{proposition}
\begin{proof}
Let $s_1=2+it$ and subtract \eqref{eq:zetapf} at $s_1$ from \eqref{eq:zetapf} at $s$. The left side at $s_1$ is bounded, by $\sum\Lambda(n)n^{-2}$. The terms $\frac1s-\frac1{s_1}$ and
$\frac1{s-1}-\frac1{s_1-1}$ are bounded since $|t|\ge2$. Both $\frac s2$ and $\frac{s_1}2$ lie in the sector $|\arg z|\le\pi-\arctan2$, with modulus at least $1$, so by Theorem~\ref{thm:stirlingc}(c),
$\psi_\Gamma(\frac s2)-\psi_\Gamma(\frac{s_1}2)=\log\frac s{s_1}+O(1)=O(1)$. There remains $\sum_\rho\bigl(\frac1{s-\rho}-\frac1{s_1-\rho}\bigr)$. If $|\gamma-t|>1$, the term is
$\frac{s_1-s}{(s-\rho)(s_1-\rho)}$, of modulus at most $\frac3{|\gamma-t|^2}\le\frac6{1+(\gamma-t)^2}$; by Corollary~\ref{cor:nunit} these add up to $O(\log|t|)$. If $|\gamma-t|\le1$, then
$|s_1-\rho|\ge2-\beta>1$, so $\bigl|\frac1{s_1-\rho}\bigr|<1$, and there are $O(\log|t|)$ such $\rho$. What remains is $\sum_{|\gamma-t|\le1}\frac1{s-\rho}$.
\end{proof}

Now comes the margin. Chapter~5 showed that no zero lies on $\operatorname{Re}s=1$. A zero could still approach the line as $t$ grows; the next theorem says how fast it may.

\begin{theorem}[de la Vallée Poussin, 1899]\label{thm:zerofree}
There is an absolute constant $c>0$ such that $\zeta(\sigma+it)\ne0$ for
\[
\sigma\ge1-\frac c{\log(|t|+2)}.
\]
\end{theorem}
\begin{proof}
\emph{Three inequalities for $1<\sigma\le2$.} First, the inequality $3+4\cos\phi+\cos2\phi\ge0$ of Theorem~\ref{thm:line}, applied with $\phi=t\log n$ and summed against $\Lambda(n)n^{-\sigma}$
(Lemma~\ref{lem:logderiv}), gives
\[
\begin{aligned}3\Bigl(-\frac{\zeta'}\zeta(\sigma)\Bigr)&+4\operatorname{Re}\Bigl(-\frac{\zeta'}\zeta(\sigma+it)\Bigr)\eqbr+\operatorname{Re}\Bigl(-\frac{\zeta'}\zeta(\sigma+2it)\Bigr)\ge0.\end{aligned}
\]
Second, $-\frac{\zeta'}\zeta(s)-\frac1{s-1}$ is holomorphic near $1$ (the proof of Theorem~\ref{thm:pnt}), hence bounded on $[1,2]$: $-\frac{\zeta'}\zeta(\sigma)\le\frac1{\sigma-1}+A$. Third, for $|t|\ge2$,
Proposition~\ref{prop:nearzeros} gives $\operatorname{Re}\bigl(-\frac{\zeta'}\zeta(\sigma+it)\bigr)=-\sum_{|\gamma-t|\le1}\operatorname{Re}\frac1{s-\rho}+O(\log|t|)$, and every term of the sum is positive: $\operatorname{Re}\frac1{s-\rho}=\frac{\sigma-\beta}{|s-\rho|^2}$
with $\sigma>1>\beta$. Dropping terms can only help: $\operatorname{Re}\bigl(-\frac{\zeta'}\zeta(\sigma+2it)\bigr)\le A\log|t|$ and, if $\rho_0=\beta_0+it$ is a zero, $\operatorname{Re}\bigl(-\frac{\zeta'}\zeta(\sigma+it)\bigr)\le-\frac1{\sigma-\beta_0}+A\log|t|$,
for a suitable constant $A\ge1$. Collecting all constants into one $A'\ge1$ (using $\log|t|\ge\log2$):

\emph{The squeeze.} Combining, for a zero $\rho_0=\beta_0+it$ with $|t|\ge2$ and every $\sigma\in\left]1,2\right]$,
\[
\frac4{\sigma-\beta_0}\le\frac3{\sigma-1}+A'\log|t|.
\]
This is the 3-4-1 argument of Chapter~5 made quantitative: the pole contributes $3$, a zero near $1+it$ would contribute $-4$, and $4>3$. Take $\sigma=1+\frac\delta{\log|t|}$ with $\delta=\frac1{2A'}$
(so $\sigma\le2$). Then $\frac4{\sigma-\beta_0}\le\bigl(\frac3\delta+A'\bigr)\log|t|=7A'\log|t|$, so $\sigma-\beta_0\ge\frac4{7A'\log|t|}=\frac{8\delta}{7\log|t|}$ and
\[
\begin{aligned}1-\beta_0&=(\sigma-\beta_0)-(\sigma-1)\eqbr\ge\frac{\delta}{7\log|t|}=\frac1{14A'\log|t|}.\end{aligned}
\]
For $|t|<2$ there are finitely many zeros with $0<\beta<1$ (they are isolated in a compact set, none on the line $\operatorname{Re}s=1$), all with $\beta<1$, and a smaller $c$ covers them.
\end{proof}

\begin{remark}
The region has been widened only slightly in more than a century: Vinogradov and Korobov (1958) reached $\sigma\ge1-\frac c{(\log|t|)^{2/3}(\log\log|t|)^{1/3}}$, which gives the best known error term in the
prime number theorem (Chapter~7). The Riemann hypothesis would give everything up to $\sigma>\frac12$. Hardy proved in 1914 that infinitely many zeros lie on the line, and Selberg in 1942 that
a positive proportion of them do. Volume~13 quotes a positive proportion and does not derive it. Numerical lower bounds on the proportion are external and have moved (past $41\%$ in recent work); they are not proved in this series.
\end{remark}

\section{Exercises}

\begin{exercise}[$\star$]\label{ex:digammahalf}
Differentiate the logarithm of Legendre's duplication formula (Volume~4, Theorem~5.14) to show that $\psi_\Gamma(\frac12)=-\gamma-2\log2$, and deduce that $B=\frac12\log4\pi-1-\frac\gamma2=-0.0230957\ldots$
\end{exercise}

\begin{exercise}[$\star\star$]\label{ex:bsum}
Show that $\sum_\rho\operatorname{Re}\frac1\rho$ converges absolutely and that $B=-\sum_\rho\operatorname{Re}\frac1\rho$. Deduce that if the Riemann hypothesis holds, then $\sum_\rho\frac1{|\rho|^2}=-2B=0.0461914\ldots$
\end{exercise}

\begin{exercise}[$\star$]\label{ex:zerosums}
Show that $\sum_\rho|\rho|^{-1}$ diverges, while $\sum_\rho|\rho|^{-1}(\log|\rho|)^{-3}$ converges.
\end{exercise}

\begin{exercise}[$\star\star$, the signature exercise]\label{ex:nthzero}
Let $0<\gamma_1\le\gamma_2\le\cdots$ be the ordinates of the zeros in the upper half-plane. Show that $\gamma_n\sim\frac{2\pi n}{\log n}$, and that the average gap between consecutive ordinates near
height $T$ is $\frac{2\pi}{\log(T/2\pi)}$. (At $T=10^6$ it is $0.524$; near the first zero the formula gives $7.8$.)
\end{exercise}

\begin{exercise}[$\star\star$]\label{ex:carlson}
Let $f$ be entire, with $|f(z)|\le Ae^{c|z|}$ for some $c<\frac{\log2}2$, and $f(n)=0$ for every integer $n\ge1$. Show that $f=0$. What does $\sin\pi z$ say about the constant?
\end{exercise}

\begin{exercise}[$\star\star$]\label{ex:zetahalf}
Show that $\zeta(s)=\sum_{n\le N}n^{-s}+\frac{N^{1-s}}{s-1}-s\int_N^\infty\{u\}u^{-s-1}\dd u$ for $\operatorname{Re}s>0$, $s\ne1$ and every integer $N\ge1$. Deduce that $\zeta(\frac12+it)=O(|t|^{1/2})$.
\end{exercise}

\begin{exercise}[$\star\star\star$]\label{ex:gammaratio}
Use \eqref{eq:stirlingvert} and the functional equation to show that $|\zeta(-1+it)|\asymp|t|^{3/2}|\zeta(2-it)|$ as $|t|\to\infty$, and more generally that
$|\zeta(\sigma+it)|=\bigl(\frac{|t|}{2\pi}\bigr)^{\frac12-\sigma}|\zeta(1-\sigma+it)|\bigl(1+O(\frac1{|t|})\bigr)$ for fixed real $\sigma$.
\end{exercise}

\section{Hints}

\begin{hint}{ex:digammahalf}
$\log\Gamma(x)+\log\Gamma(x+\frac12)=(1-2x)\log2+\frac12\log\pi+\log\Gamma(2x)$; take $x=\frac12$ after differentiating, and use $\psi_\Gamma(1)=\Gamma'(1)=-\gamma$. Then compute $\frac{\xi'}\xi(1)=-B$.
\end{hint}
\begin{hint}{ex:bsum}
Use Theorem~\ref{thm:hadamard} at $s=1$, $\frac{\xi'}\xi(1)=-B$, and the symmetry $\rho\mapsto1-\overline\rho$ of the zeros. Under the Riemann hypothesis $\operatorname{Re}\frac1\rho=\frac1{2|\rho|^2}$.
\end{hint}
\begin{hint}{ex:zerosums}
Count the zeros with $2^k<|\rho|\le2^{k+1}$ with Theorem~\ref{thm:rvm}.
\end{hint}
\begin{hint}{ex:nthzero}
$N(\gamma_n)\approx n$; invert $\frac T{2\pi}\log\frac T{2\pi}$ as in Exercise~\ref{ex:nthprime}.
\end{hint}
\begin{hint}{ex:carlson}
Apply Corollary~\ref{cor:jensen} to $f(z+1)$ or to $f$ on a disc centred at a point where $f\ne0$; the disc of radius $r$ contains about $r$ integers.
\end{hint}
\begin{hint}{ex:zetahalf}
Summation by parts on $\sum_{N<n\le X}n^{-s}$ for $\operatorname{Re}s>1$, then analytic continuation. Take $N=\lfloor|t|\rfloor$.
\end{hint}
\begin{hint}{ex:gammaratio}
Write the functional equation as
\[\zeta(s)=\pi^{s-\frac12}\frac{\Gamma((1-s)/2)}{\Gamma(s/2)}\zeta(1-s).\]
Compare $|\Gamma(\frac{1-\sigma}2-\frac{it}2)|$ and $|\Gamma(\frac\sigma2+\frac{it}2)|$ with \eqref{eq:stirlingvert}.
\end{hint}

\section{Solutions}

\begin{solution}{ex:digammahalf}
\idea A known value of $\psi_\Gamma$ comes from an identity for $\Gamma$; then $B$ is $-\frac{\xi'}\xi(1)$, computed factor by factor.
\writeup Differentiating the logarithm of the duplication formula, $\psi_\Gamma(x)+\psi_\Gamma(x+\frac12)=-2\log2+2\psi_\Gamma(2x)$. At $x=\frac12$: $\psi_\Gamma(\frac12)+\psi_\Gamma(1)=-2\log2+2\psi_\Gamma(1)$, so
$\psi_\Gamma(\frac12)=\psi_\Gamma(1)-2\log2=-\gamma-2\log2$ (Volume~4, Exercise~5.7 gives $\Gamma'(1)=-\gamma$).

From $\xi(s)=\xi(1-s)$, $\xi'(0)=-\xi'(1)$, so $B=-\frac{\xi'}\xi(1)$. Near $s=1$, $\frac{\xi'}\xi(s)=\frac1s+\bigl(\frac1{s-1}+\frac{\zeta'}\zeta(s)\bigr)-\frac12\log\pi+\frac12\psi_\Gamma(\frac s2)$, and
$\frac1{s-1}+\frac{\zeta'}\zeta(s)\to\gamma$ (in the notation of Chapter~5, $\frac{\zeta'}\zeta=\frac{Z'}Z-\frac1{s-1}$ and $Z'(1)=\gamma$). So
\[
\begin{aligned}\frac{\xi'}\xi(1)&=1+\gamma-\frac12\log\pi-\frac\gamma2-\log2\\&=1+\frac\gamma2-\frac12\log4\pi,\end{aligned}
\]
and $B=\frac12\log4\pi-1-\frac\gamma2=1.2655121-1-0.2886078=-0.0230957$.
\end{solution}

\begin{solution}{ex:bsum}
\idea The partial fractions at $s=1$ involve $\frac1{1-\rho}+\frac1\rho$; the symmetry of the zeros turns $\frac1{1-\rho}$ into another $\frac1{\overline{\rho'}}$, whose real part is that of $\frac1{\rho'}$.
\writeup $\operatorname{Re}\frac1\rho=\frac\beta{|\rho|^2}\in\left]0,\frac1{|\rho|^2}\right[$, so the series converges absolutely. Theorem~\ref{thm:hadamard} at $s=1$ gives $-B=\frac{\xi'}\xi(1)=B+\sum_\rho\bigl(\frac1{1-\rho}+\frac1\rho\bigr)$,
the series converging absolutely. Take real parts. The map $\rho\mapsto\rho^*=1-\overline\rho$ permutes the zeros (Proposition~\ref{prop:xi}(b)) and $1-\rho=\overline{\rho^*}$, so
$\operatorname{Re}\frac1{1-\rho}=\operatorname{Re}\frac1{\rho^*}$; summing a convergent series of positive terms in another order does not change it, so $\sum_\rho\operatorname{Re}\frac1{1-\rho}=\sum_\rho\operatorname{Re}\frac1\rho$. As $B$ is real,
$-2B=2\sum_\rho\operatorname{Re}\frac1\rho$. Under the Riemann hypothesis $\beta=\frac12$, so $\sum|\rho|^{-2}=2\sum\operatorname{Re}\frac1\rho=-2B=0.0461914$.
\end{solution}

\begin{solution}{ex:zerosums}
\idea There are about $\frac T{2\pi}\log T$ zeros up to height $T$, so each dyadic block $2^k<|\rho|\le2^{k+1}$ holds about $k2^k$ of them, each of size $2^k$.
\writeup Since $|\rho|$ and $|\gamma|$ differ by less than $1$ and zeros come in conjugate pairs, the number of zeros with $2^k<|\rho|\le2^{k+1}$ is $2\bigl(N(2^{k+1})-N(2^k)\bigr)+O(\log2^k)$, which by
Theorem~\ref{thm:rvm} is $\frac{2^k}\pi\bigl(k\log2+O(1)\bigr)$, between $c_1k2^k$ and $c_2k2^k$ for $k$ large. In this block $|\rho|^{-1}\ge2^{-k-1}$, so the block contributes at least $\frac{c_1k}2$ to $\sum|\rho|^{-1}$:
divergence. With the extra factor $(\log|\rho|)^{-3}\le(k\log2)^{-3}$ the block contributes $O(k^{-2})$, a convergent series.
\end{solution}

\begin{solution}{ex:nthzero}
\idea $N(T)\approx\frac T{2\pi}\log T$, so inverting gives $T\approx\frac{2\pi n}{\log n}$, exactly as $\pi(x)\approx\frac x{\log x}$ inverted to $p_n\approx n\log n$, with the roles of multiplying and
dividing by $\log$ exchanged.
\writeup $N(\gamma_n)\ge n$ and $N(\gamma_n-\varepsilon)<n$ for small $\varepsilon$, so $N(\gamma_n)=n+O(\log\gamma_n)$, and Theorem~\ref{thm:rvm} gives $\frac{\gamma_n}{2\pi}\log\gamma_n\sim n$. Taking logarithms,
$\log\gamma_n\sim\log n$, so $\gamma_n\sim\frac{2\pi n}{\log n}$. The number of ordinates in $]T,T+H]$ for $H\to\infty$, $H=o(T)$, is $\frac H{2\pi}\log\frac T{2\pi}+O(\log T)$, since the derivative of the main term is
$\frac1{2\pi}\log\frac T{2\pi}$; so the average gap is $\frac{2\pi}{\log(T/2\pi)}$. At $T=10^6$, $\log\frac{10^6}{2\pi}=11.98$ and the gap is $0.524$.
Near the first zero, $\log\frac{14}{2\pi}=0.80$ and the formula gives $7.8$, compared with the actual gaps $6.9$, $4.0$, $5.4$ between the first four ordinates.
\end{solution}

\begin{solution}{ex:carlson}
\idea Zeros are expensive for a slowly growing function (Jensen). Integers are evenly spread, about $r$ of them in a disc of radius $r$, while the budget is about $2cr$ for a disc of
radius $2r$.
\writeup Suppose $f\ne0$. Replacing $f$ by $\frac{f(z)}{z^m}$ for a suitable $m$ changes neither the zeros at the positive integers nor the growth bound (up to the constant $A$, on $|z|\ge1$), so
assume $f(0)\ne0$. The disc $|z|\le r$ contains the $\lfloor r\rfloor$ zeros $1,\dots,\lfloor r\rfloor$. Corollary~\ref{cor:jensen} with $R=2r$ gives $\lfloor r\rfloor\log2\le\log\frac{Ae^{2cr}}{|f(0)|}=2cr+O(1)$. Dividing
by $r$ and letting $r\to\infty$: $\log2\le2c$, a contradiction. (Optimising the ratio $\frac Rr$ improves the constant; Carlson's theorem gives the sharp result, with growth $e^{c|z|}$ for any $c<\pi$ along the
imaginary axis. The function $\sin\pi z$ vanishes at all integers and has $|\sin\pi z|\le e^{\pi|z|}$, so some bound of this type is necessary.)
\end{solution}

\begin{solution}{ex:zetahalf}
\idea Summation by parts writes the tail of $\sum n^{-s}$ as an integral minus a fluctuating term; the formula then extends to $\operatorname{Re}s>0$, and cutting at $N\approx|t|$ balances the two parts.
\writeup For $\operatorname{Re}s>1$ and $X>N$, Lemma~\ref{lem:abelint} with $a_n=1$ for $n>N$, $A(u)=\lfloor u\rfloor-N$, $f(u)=u^{-s}$, gives
$\sum_{N<n\le X}n^{-s}=(\lfloor X\rfloor-N)X^{-s}+s\int_N^X(\lfloor u\rfloor-N)u^{-s-1}\dd u$. Write $\lfloor u\rfloor=u-\{u\}$ and let $X\to\infty$:
$\sum_{n>N}n^{-s}=s\int_N^\infty u^{-s}\dd u-Ns\int_N^\infty u^{-s-1}\dd u-s\int_N^\infty\{u\}u^{-s-1}\dd u=\frac{N^{1-s}}{s-1}-s\int_N^\infty\{u\}u^{-s-1}\dd u$. The right side is holomorphic on
$\operatorname{Re}s>0$ except at $1$ (the integral is bounded by $\int_N^\infty u^{-\sigma-1}$ uniformly on compact sets; split it at the integers to apply Volume~4, Theorem~9.15), so the identity extends.

At $s=\frac12+it$, $|t|\ge2$, $N=\lfloor|t|\rfloor$: $\bigl|\sum_{n\le N}n^{-1/2}\bigr|\le2\sqrt N$, $\bigl|\frac{N^{1/2-it}}{s-1}\bigr|\le\frac{\sqrt N}{|t|}\le1$, and $\bigl|s\int_N^\infty\{u\}u^{-3/2-it}\dd u\bigr|\le|s|\cdot\frac2{\sqrt N}$, which is $O(\sqrt{|t|})$.
So $\zeta(\frac12+it)=O(|t|^{1/2})$. This is far from the truth: Hardy and Littlewood reached $|t|^{1/4+\varepsilon}$ and Weyl $|t|^{1/6+\varepsilon}$, Bourgain (2017) $|t|^{13/84+\varepsilon}$, and the Lindelöf
hypothesis, a consequence of the Riemann hypothesis, predicts $|t|^\varepsilon$.
\end{solution}

\begin{solution}{ex:gammaratio}
\idea The functional equation relates $\zeta$ at $s$ and $1-s$ through a ratio of Gamma values; on vertical lines Stirling says that the exponentials $e^{-\pi|t|/4}$ cancel and a power of $|t|$
survives.
\writeup From $\xi(s)=\xi(1-s)$ we get $\pi^{-s/2}\Gamma(\frac s2)\zeta(s)=\pi^{-(1-s)/2}\Gamma(\frac{1-s}2)\zeta(1-s)$, which is the formula of the hint. Let $s=\sigma+it$, $t\to+\infty$. By \eqref{eq:stirlingvert},
$\log|\Gamma(a+ib)|=(a-\frac12)\log b-\frac{\pi b}2+\frac12\log2\pi+O(\frac1b)$ for fixed $a$; also $|\Gamma(\overline z)|=|\Gamma(z)|$. With $b=\frac t2$:
\[
\begin{aligned}\log\Bigl|\frac{\Gamma(\frac{1-\sigma}2-\frac{it}2)}{\Gamma(\frac\sigma2+\frac{it}2)}\Bigr|&=\Bigl(\frac{1-\sigma}2-\frac\sigma2\Bigr)\log\frac t2+O\Bigl(\frac1t\Bigr)\\&=\Bigl(\frac12-\sigma\Bigr)\log\frac t2+O\Bigl(\frac1t\Bigr).\end{aligned}
\]
And $|\pi^{s-1/2}|=\pi^{\sigma-1/2}$. So $|\zeta(\sigma+it)|=\bigl(\frac t{2\pi}\bigr)^{\frac12-\sigma}|\zeta(1-\sigma-it)|(1+O(\frac1t))$, and $|\zeta(1-\sigma-it)|=|\zeta(1-\sigma+it)|$ by conjugation. For $\sigma=-1$:
$|\zeta(-1+it)|\sim\bigl(\frac t{2\pi}\bigr)^{3/2}|\zeta(2+it)|$, and $|\zeta(2+it)|$ lies between $\frac{\zeta(4)}{\zeta(2)}$ and $\zeta(2)$. In particular, on the left of the strip $\zeta$ grows like a power of $|t|$.
\end{solution}

\section*{Sources and further reading}

\begin{itemize}
\item B.~Riemann, \emph{Ueber die Anzahl der Primzahlen unter einer gegebenen Grösse}, Monatsberichte der Berliner Akademie (1859), 671--680. English translation in H.~M.~Edwards's book below.
\item H.~M.~Edwards, \emph{Riemann's Zeta Function}, Academic Press, 1974 (Dover reprint). A line-by-line commentary on Riemann's memoir and what became of each claim.
\item J.~Hadamard, \emph{Étude sur les propriétés des fonctions entières et en particulier d'une fonction considérée par Riemann}, Journal de Mathématiques Pures et Appliquées 9 (1893), 171--215.
\item J.~L.~W.~V.~Jensen, \emph{Sur un nouvel et important théorème de la théorie des fonctions}, Acta Mathematica 22 (1899), 359--364.
\item H.~Davenport, \emph{Multiplicative Number Theory}, 3rd ed., Springer GTM~74, 2000, Chapters~11--15 (Hadamard's product, the zero-free region, $N(T)$).
\item H.~L.~Montgomery and R.~C.~Vaughan, \emph{Multiplicative Number Theory I}, Cambridge, 2007, Chapters~6 and~10 (the local lemma in the form used here, Lemma~6.3).
\item E.~C.~Titchmarsh, \emph{The Theory of the Riemann Zeta-Function}, 2nd ed., revised by D.~R.~Heath-Brown, Oxford, 1986, Chapters~9 and~10.
\item A.~M.~Turing, \emph{Some calculations of the Riemann zeta-function}, Proceedings of the London Mathematical Society 3 (1953), 99--117.
\end{itemize}
